\chapter{General Project Context}

\section{Introduction}
This chapter introduces the academic and industrial context of \projectname. It clarifies the company environment, the broader strategic initiative, the motivation for the project, and the positioning of the PFE work within the larger transformation objective.

\section{Company and Strategic Context}
\projectname{} is developed within \initiative{} at \companyname. The global initiative aims to strengthen data-driven marketing and customer intelligence capabilities. Within this larger vision, the PFE contribution focuses specifically on the distributed and secure data platform required to centralize, process, and expose customer data reliably.

\section{Problem Statement}
Traditional customer data environments are often fragmented across multiple channels, technologies, and business silos. This fragmentation reduces data consistency, limits analytical agility, and makes future intelligent services more difficult to implement. A distributed lakehouse platform addresses these issues by combining scalable storage, structured governance, analytical query access, and orchestrated processing.

\section{Project Positioning}
The project must be defended as a distributed customer data engineering platform rather than as a standalone dashboard, a simple ETL chain, or a finished AI model. Its purpose is to provide the infrastructure and governance foundation upon which future Customer 360 and machine learning services can be built.

\section{Chapter Conclusion}
This chapter establishes the academic and business framing of the project and prepares the transition toward the theoretical foundations presented in the next chapter.

